% !TeX root = draft_finalized.tex
% Appendix for the identification revision. Formatting matches body.tex
% (booktabs + threeparttable + L{} columns, elsarticle-harv citations).

\appendix
\renewcommand{\thesection}{\Alph{section}}
\renewcommand{\thetable}{A\arabic{table}}
\renewcommand{\thefigure}{A\arabic{figure}}
\setcounter{table}{0}
\setcounter{figure}{0}

\section{Supplementary tables}
\label{app:robust}

This appendix reports the sample composition, treatment distribution, and the additional
identification and inference checks summarized in Section~\ref{sec:results}. Unless noted,
all specifications use the baseline design: survey-weighted linear probability models with
historical-province, birth-cohort, and survey-year fixed effects, controls for gender (pooled)
and urban residence, and standard errors two-way clustered by historical province and birth
cohort. Vertical and joint outcomes use the unrestricted full-period occupational education
benchmark; as Table~\ref{tab:robdef} shows, the preferred-benchmark point estimates differ by
less than 0.01 percentage points.

\begin{table}[htbp]
  \centering
  \caption{Analytical samples and outcome means}
  \label{tab:appsamples}
  \begin{threeparttable}
  \begin{tabular}{L{6.4cm}rr}
    \toprule
    Outcome / sample & Observations & Weighted mean \\
    \midrule
    Broad analysis sample (exposure, cohort, geography) & 4{,}086{,}861 & --- \\
    \addlinespace
    Labor-force participation      & 4{,}086{,}861 & 0.496 \\
    Employment                     & 4{,}086{,}861 & 0.421 \\
    Unemployment (participants)     & 1{,}996{,}995 & 0.151 \\
    \addlinespace
    Overeducation                  & 1{,}633{,}647 & 0.361 \\
    Undereducation                 & 1{,}633{,}647 & 0.238 \\
    Any vertical mismatch          & 1{,}633{,}647 & 0.599 \\
    \addlinespace
    Horizontal mismatch (broad)    & 400{,}726 & 0.510 \\
    Double mismatch                & 400{,}726 & 0.366 \\
    Fully matched employment       & 400{,}726 & 0.235 \\
    \bottomrule
  \end{tabular}
  \begin{tablenotes}[flushleft]
    \footnotesize
    \item Observation counts are for the pooled sample. Weighted means use LFS survey weights.
    The vertical-mismatch sample requires employment and a valid occupation code; the horizontal
    and joint samples additionally require tertiary education, a valid field of study, and a
    harmonized two-digit ISCO code.
  \end{tablenotes}
  \end{threeparttable}
\end{table}

\begin{table}[htbp]
  \centering
  \caption{University-exposure distribution and variance decomposition}
  \label{tab:appexposure}
  \begin{threeparttable}
  \begin{tabular}{L{7.4cm}r}
    \toprule
    \multicolumn{2}{l}{\textit{Panel A. Exposure per 100{,}000 population (ages 18--20 average)}} \\
    \addlinespace
    Weighted mean & 739.7 \\
    Weighted standard deviation & 535.7 \\
    10th / 25th / 50th / 75th / 90th percentile & 194 / 292 / 582 / 1197 / 1490 \\
    Minimum / maximum & 107 / 4702 \\
    \addlinespace
    \multicolumn{2}{l}{\textit{Panel B. Log exposure (identifying variable)}} \\
    \addlinespace
    Raw weighted standard deviation & 0.76 \\
    Raw weighted interquartile range & 1.32 \\
    Residualized standard deviation (after prov.\ $+$ cohort $+$ year FE) & 0.14 \\
    Residualized interquartile range & 0.19 \\
    \addlinespace
    \multicolumn{2}{l}{\textit{Panel C. Variance decomposition of log exposure}} \\
    \addlinespace
    Province--cohort cells & 804 \\
    Share absorbed by province $+$ cohort $+$ survey-year fixed effects & 0.96 \\
    Residual (identifying) share & 0.04 \\
    \bottomrule
  \end{tabular}
  \begin{tablenotes}[flushleft]
    \footnotesize
    \item Panel A describes the individual-level exposure measure on the broad analysis sample.
    Panels B and C describe the natural logarithm of exposure, which enters the regressions.
    The residualized moments in Panel B are computed from the residuals of a survey-weighted
    regression of log exposure on the identifying fixed effects on the joint-mismatch sample, and
    give the scale of the variation that actually identifies the coefficient. Panel C reports the
    share of the weighted variance of log exposure absorbed by those fixed effects.
  \end{tablenotes}
  \end{threeparttable}
\end{table}

\begin{table}[htbp]
  \centering
  \caption{University exposure and educational attainment}
  \label{tab:appfirststage}
  \begin{threeparttable}
  \begin{tabular}{L{7.6cm}cc}
    \toprule
    Specification & Effect (pp) & $p$-value \\
    \midrule
    Tertiary education, province FE only                 & $0.79^{***}$ $(0.07)$ & $<0.001$ \\
    Tertiary education, $+$ birth-cohort FE              & $-0.15$ $(0.17)$ & 0.36 \\
    Tertiary education, baseline FE (prov $+$ cohort $+$ year) & $-0.15$ $(0.17)$ & 0.36 \\
    Tertiary education, $+$ province $\times$ year FE     & $-0.16$ $(0.17)$ & 0.33 \\
    Tertiary education, $+$ province-specific cohort trend & $-0.03$ $(0.19)$ & 0.87 \\
    \addlinespace
    Tertiary education, ages $\geq 25$                    & $-0.05$ $(0.14)$ & 0.71 \\
    Tertiary education, ages $\geq 30$                    & $-0.03$ $(0.13)$ & 0.78 \\
    Tertiary education, women, ages $\geq 25$             & $-0.10$ $(0.14)$ & 0.45 \\
    Tertiary education, men, ages $\geq 25$               & $-0.00$ $(0.16)$ & 0.99 \\
    \addlinespace
    Education attainment rank (1--8), ages $\geq 25$      & $0.006$ $(0.005)$ & 0.24 \\
    Literacy, ages $\geq 25$                              & $0.014$ $(0.015)$ & 0.34 \\
    \bottomrule
  \end{tabular}
  \begin{tablenotes}[flushleft]
    \footnotesize
    \item Each row is a separate regression of the indicated educational outcome on log exposure.
    Effects are percentage-point changes associated with a 10 percent increase in exposure
    (attainment-rank effect is in rank points). The province-FE-only row identifies the
    coefficient from cross-province variation and is reported only to show that the cross-sectional
    association vanishes once cohort variation is isolated. Residualized-exposure quartiles of the
    ages-$\geq 25$ sample have tertiary shares of 0.31, 0.26, 0.27, and 0.29 from lowest to highest.
    \textsuperscript{***}$p<0.01$.
  \end{tablenotes}
  \end{threeparttable}
\end{table}

\begin{table}[htbp]
  \centering
  \caption{Population-level (unconditional) mismatch incidence}
  \label{tab:appuncond}
  \begin{threeparttable}
  \begin{tabular}{L{6.8cm}ccc}
    \toprule
    Outcome (share of broad population) & Baseline FE & \makecell{$+$ Province $\times$\\year FE} & Mean \\
    \midrule
    Tertiary-educated and employed                    & $0.015$ $(0.050)$ & $0.012$ $(0.048)$ & 0.117 \\
    Tertiary-educated, employed, horizontally mismatched & $0.044^{*}$ $(0.026)$ & $0.039$ $(0.026)$ & 0.059 \\
    Tertiary-educated, employed, doubly mismatched     & $0.041$ $(0.028)$ & $0.038$ $(0.026)$ & 0.043 \\
    Tertiary-educated, employed, fully matched         & $-0.011$ $(0.018)$ & $-0.011$ $(0.018)$ & 0.027 \\
    \bottomrule
  \end{tabular}
  \begin{tablenotes}[flushleft]
    \footnotesize
    \item The outcome is an indicator defined over the entire broad analysis sample (not
    conditional on tertiary education or employment), equal to one for individuals observed in
    the indicated tertiary-educated employed state and zero otherwise. Effects are percentage-point
    changes for a 10 percent increase in exposure. These estimates bound the population-level
    contribution of exposure to graduate mismatch and are smaller than the conditional estimates
    in Table~\ref{tab:mismatch}, consistent with the null first stage in
    Table~\ref{tab:appfirststage}. \textsuperscript{*}$p<0.10$.
  \end{tablenotes}
  \end{threeparttable}
\end{table}

\begin{table}[htbp]
  \centering
  \caption{Frozen occupational education benchmark}
  \label{tab:appfrozen}
  \begin{threeparttable}
  \begin{tabular}{L{7.0cm}cc}
    \toprule
    Outcome & Effect (pp) & Observations \\
    \midrule
    Overeducation (frozen benchmark)                       & $0.366^{***}$ $(0.120)$ & 1{,}523{,}305 \\
    Undereducation (frozen benchmark)                      & $-0.077$ $(0.051)$ & 1{,}523{,}305 \\
    Double mismatch (frozen vertical $+$ broad horizontal) & $0.705^{***}$ $(0.145)$ & 343{,}371 \\
    Fully matched employment (frozen benchmark)            & $-0.298^{**}$ $(0.110)$ & 343{,}371 \\
    \bottomrule
  \end{tabular}
  \begin{tablenotes}[flushleft]
    \footnotesize
    \item The frozen benchmark is the survey-weighted modal education rank within each two-digit
    ISCO occupation computed using only the earliest six survey years (1386--1391) and applied to
    all observations; occupations with fewer than 200 weighted observations in the benchmark
    window are dropped. This benchmark cannot be mechanically influenced by university expansion
    in later years. Effects are percentage-point changes for a 10 percent increase in exposure.
    \textsuperscript{***}$p<0.01$, \textsuperscript{**}$p<0.05$.
  \end{tablenotes}
  \end{threeparttable}
\end{table}

\begin{table}[htbp]
  \centering
  \caption{Inference with a limited number of clusters}
  \label{tab:appcluster}
  \begin{threeparttable}
  \begin{tabular}{L{4.6cm}cccc}
    \toprule
    Outcome (pooled) & \makecell{Two-way\\cluster $p$} & \makecell{Province\\cluster $p$} & \makecell{Cohort\\cluster $p$} & \makecell{Leave-one-province-out\\effect range (pp)} \\
    \midrule
    Horizontal mismatch      & $<0.001$ & $<0.001$ & $<0.001$ & $0.52$ to $0.66$ \\
    Double mismatch          & $<0.001$ & $<0.001$ & $<0.001$ & $0.62$ to $0.73$ \\
    Fully matched employment & 0.003 & 0.002 & $<0.001$ & $-0.40$ to $-0.28$ \\
    Overeducation            & 0.003 & 0.007 & 0.008 & --- \\
    \bottomrule
  \end{tabular}
  \begin{tablenotes}[flushleft]
    \footnotesize
    \item Columns 2--4 report $p$-values for the coefficient on log exposure under two-way
    clustering (baseline) and one-way clustering by historical province (28 clusters) and by
    birth cohort (29 clusters); source \texttt{step16\_main\_inference\_results.csv}. Column 5
    reports the range of the estimated effect across the 28 samples that each omit one historical
    province; in every leave-one-out sample the horizontal- and double-mismatch coefficients
    retain their sign and remain significant at the 5 percent level, and the fully-matched
    coefficient is significant at 5 percent in 27 of 28 samples (source
    \texttt{step16\_leave\_one\_province\_out\_summary.csv}). A restricted wild cluster bootstrap
    (Rademacher weights, imposed null) was also implemented but did not converge within available
    computing resources on the full samples; the one-way and leave-one-out results above serve as
    the few-cluster check.
  \end{tablenotes}
  \end{threeparttable}
\end{table}

\begin{table}[htbp]
  \centering
  \caption{Formal tests of gender differences}
  \label{tab:appgender}
  \begin{threeparttable}
  \begin{tabular}{L{4.6cm}cccc}
    \toprule
    Outcome & \makecell{Male\\effect (pp)} & \makecell{Female\\effect (pp)} & \makecell{Female $-$ male\\difference (pp)} & \makecell{$p$-value\\(difference)} \\
    \midrule
    Overeducation            & $0.300$   & $0.479$  & $0.179$ $(0.213)$  & 0.41 \\
    Undereducation           & $-0.157$  & $-0.181$ & $-0.024$ $(0.144)$ & 0.87 \\
    Any vertical mismatch    & $0.142$   & $0.297$  & $0.155$ $(0.192)$  & 0.43 \\
    Horizontal mismatch      & $0.840$   & $0.073$  & $-0.767$ $(0.213)$ & $0.001$ \\
    Double mismatch          & $0.784$   & $0.367$  & $-0.417$ $(0.236)$ & 0.089 \\
    Fully matched employment & $-0.519$  & $0.094$  & $0.613$ $(0.245)$  & 0.019 \\
    Employment               & $-0.052$  & $-0.167$ & $-0.115$ $(0.268)$ & 0.67 \\
    Labor-force participation & $0.078$  & $-0.161$ & $-0.239$ $(0.255)$ & 0.36 \\
    \bottomrule
  \end{tabular}
  \begin{tablenotes}[flushleft]
    \footnotesize
    \item Each row is a single pooled regression of the outcome on log exposure, a female
    indicator, their interaction, urban residence, and urban $\times$ female, with
    \emph{gender-interacted} historical-province, birth-cohort, and survey-year fixed effects
    ($\text{province}\times\text{female}$, $\text{cohort}\times\text{female}$,
    $\text{year}\times\text{female}$) and standard errors two-way clustered by province and cohort.
    This specification reproduces the separate-sex regressions exactly, so the coefficient on
    log exposure $\times$ female is a clean test of whether the male and female exposure
    coefficients differ. Effects are percentage-point changes for a 10 percent increase in
    exposure. Under this test the male--female difference is statistically significant for
    horizontal mismatch and fully matched employment, significant only at the 10 percent level for
    double mismatch, and not significant for the vertical-mismatch or employment outcomes.
  \end{tablenotes}
  \end{threeparttable}
\end{table}

\begin{table}[htbp]
  \centering
  \caption{Province-by-cohort-by-year cell-level estimates}
  \label{tab:appcell}
  \begin{threeparttable}
  \begin{tabular}{L{5.6cm}cc}
    \toprule
    Outcome & Cell-level effect (pp) & Individual-level (Table~\ref{tab:mismatch}) \\
    \midrule
    Labor-force participation & $-0.019$ $(0.061)$ & $-0.035$ $(0.057)$ \\
    Employment                & $-0.081$ $(0.086)$ & $-0.103$ $(0.085)$ \\
    Unemployment              & $0.242$ $(0.144)$  & $0.228$ $(0.147)$ \\
    Overeducation             & $0.339^{***}$ $(0.101)$ & $0.330^{***}$ $(0.101)$ \\
    Undereducation            & $-0.169^{*}$ $(0.086)$  & $-0.159^{*}$ $(0.093)$ \\
    Horizontal mismatch       & $0.590^{***}$ $(0.116)$ & $0.614^{***}$ $(0.126)$ \\
    Double mismatch           & $0.639^{***}$ $(0.136)$ & $0.688^{***}$ $(0.140)$ \\
    Fully matched employment  & $-0.315^{***}$ $(0.098)$ & $-0.346^{***}$ $(0.102)$ \\
    \bottomrule
  \end{tabular}
  \begin{tablenotes}[flushleft]
    \footnotesize
    \item The left column collapses each outcome to survey-weighted means within
    province $\times$ birth-cohort $\times$ survey-year cells (about 11{,}000 cells) and estimates
    the baseline specification by weighted least squares with cell population weights and the same
    fixed effects and clustering. The similarity to the individual-level estimates indicates that
    precision is not an artefact of the number of individual observations.
    \textsuperscript{***}$p<0.01$, \textsuperscript{*}$p<0.10$.
  \end{tablenotes}
  \end{threeparttable}
\end{table}

\begin{table}[htbp]
  \centering
  \caption{Field-of-study heterogeneity: sample sizes, adjusted $p$-values, and differences from education}
  \label{tab:appfield}
  \begin{threeparttable}
  \begin{tabular}{L{4.6cm}rccc}
    \toprule
    Field of study & \makecell{$N$\\(horizontal)} & \makecell{FDR $p$\\horizontal} & \makecell{FDR $p$\\double} & \makecell{Difference from\\education (horiz.\,/\,double)} \\
    \midrule
    Education                                  & 22{,}738  & 0.60 & 0.56 & --- \\
    Arts and humanities                        & 38{,}599  & 0.60 & 0.24 & $1.31$ / $0.95$ \\
    Social sciences and journalism             & 24{,}033  & 0.37 & 0.24 & $1.02$ / $0.94$ \\
    Business, administration, and law          & 102{,}258 & $0.009^{***}$ & $0.008^{***}$ & $1.31^{***}$ / $1.47^{***}$ \\
    Natural sciences and mathematics           & 26{,}316  & 0.37 & 0.39 & $1.01^{*}$ / $0.58$ \\
    Information and communication technologies  & 17{,}281  & $0.10^{*}$ & 0.24 & $1.52^{**}$ / $1.02^{**}$ \\
    Engineering, manufacturing, construction   & 113{,}344 & 0.37 & 0.24 & $0.75$ / $0.80$ \\
    Agriculture and veterinary studies         & 15{,}833  & 0.37 & 0.24 & $0.82^{**}$ / $1.08^{**}$ \\
    Health and welfare                         & 26{,}457  & 0.68 & 0.96 & $0.51$ / $0.31$ \\
    Services                                   & 13{,}145  & 0.37 & 0.39 & $0.70$ / $0.94$ \\
    \bottomrule
  \end{tabular}
  \begin{tablenotes}[flushleft]
    \footnotesize
    \item Point estimates and unadjusted significance are in Table~\ref{tab:field}. FDR $p$-values
    are Benjamini--Hochberg adjustments within outcome across the ten primary pooled field tests.
    ``Difference from education'' is the field-specific effect minus the effect for education
    graduates (percentage points, 10 percent exposure increase), with significance from the
    corresponding interaction test. Only business, administration, and law differs significantly
    from education after adjustment on both horizontal and double mismatch. Source:
    \texttt{step15B\_main\_pooled\_primary\_results.csv},
    \texttt{step15B\_formal\_differences\_from\_education.csv}. ICT is observed over seven survey
    years. \textsuperscript{***}$p<0.01$, \textsuperscript{**}$p<0.05$, \textsuperscript{*}$p<0.10$.
  \end{tablenotes}
  \end{threeparttable}
\end{table}

\begin{table}[htbp]
  \centering
  \caption{Robustness to age restrictions (post-treatment samples)}
  \label{tab:appage}
  \begin{threeparttable}
  \begin{tabular}{L{4.8cm}cccc}
    \toprule
    Outcome & Baseline & Age $\geq 21$ & Age $\geq 25$ & Age $\geq 30$ \\
    \midrule
    P(tertiary education)     & $-0.154$ $(0.165)$ & $-0.090$ $(0.156)$ & $-0.051$ $(0.137)$ & $-0.034$ $(0.120)$ \\
    Labor-force participation & $-0.035$ $(0.057)$ & $-0.022$ $(0.053)$ & $-0.047$ $(0.060)$ & $-0.007$ $(0.075)$ \\
    Employment                & $-0.103$ $(0.085)$ & $-0.111$ $(0.084)$ & $-0.147$ $(0.088)$ & $-0.090$ $(0.089)$ \\
    Unemployment              & $0.228$ $(0.147)$  & $0.231$ $(0.149)$  & $0.229$ $(0.141)$  & $0.159^{*}$ $(0.090)$ \\
    Overeducation             & $0.330^{***}$ $(0.101)$ & $0.309^{***}$ $(0.094)$ & $0.249^{**}$ $(0.091)$ & $0.220^{**}$ $(0.088)$ \\
    Undereducation            & $-0.159^{*}$ $(0.093)$ & $-0.119$ $(0.104)$ & $-0.081$ $(0.110)$ & $-0.046$ $(0.122)$ \\
    Horizontal mismatch       & $0.614^{***}$ $(0.126)$ & $0.605^{***}$ $(0.130)$ & $0.621^{***}$ $(0.135)$ & $0.590^{***}$ $(0.117)$ \\
    Double mismatch           & $0.688^{***}$ $(0.140)$ & $0.680^{***}$ $(0.145)$ & $0.679^{***}$ $(0.149)$ & $0.684^{***}$ $(0.164)$ \\
    Fully matched employment  & $-0.346^{***}$ $(0.102)$ & $-0.345^{***}$ $(0.103)$ & $-0.331^{***}$ $(0.105)$ & $-0.346^{***}$ $(0.113)$ \\
    \bottomrule
  \end{tabular}
  \begin{tablenotes}[flushleft]
    \footnotesize
    \item Each cell is a separate baseline regression restricted to workers at least the indicated
    age at the time of the survey. Effects are percentage-point changes for a 10 percent increase
    in exposure, with two-way clustered standard errors in parentheses. Because exposure is
    measured at ages 18--20, the age-restricted samples ensure that outcomes are observed after
    the treatment window. The horizontal-mismatch, double-mismatch, and fully-matched-employment
    estimates are essentially unchanged; overeducation attenuates but remains significant.
    \textsuperscript{***}$p<0.01$, \textsuperscript{**}$p<0.05$, \textsuperscript{*}$p<0.10$.
  \end{tablenotes}
  \end{threeparttable}
\end{table}

\subsection*{Reproducibility statement}

The individual-level analysis file is built from the Iranian Labor Force Survey public-use
microdata (Persian calendar years 1386--1403) with a set of construction scripts that harmonize
provinces to 28 historical units, harmonize education to eight ordered ranks, harmonize occupations
to two-digit ISCO and fields of study to broad ISCED groups, merge the province-year
university-enrollment and population panel, and build the vertical, horizontal, and joint mismatch
measures. All regressions are survey-weighted linear probability models estimated in \texttt{R}
(\texttt{data.table}, \texttt{fixest}); the manuscript's headline estimates were independently
reproduced from the analysis file. Every table and figure is generated by code. The field-to-occupation
crosswalks are author-constructed (mapping broad ISCED fields to two-digit ISCO groups) and are
released with the replication materials. The construction scripts, regression scripts, crosswalks,
and the province-year exposure panel are available from the corresponding author; the underlying
Labor Force Survey microdata are restricted and must be obtained from the Statistical Center of Iran.

\clearpage
